\chapter{The validation suite}
\label{sec:validation}

A tuning product can fool its builders the same way a configuration
fools a tuner: by being evaluated on the problems it was tuned on.
This document has met the trap twice, as seed overfitting in
Chapter~\ref{sec:workloads} and as the proxy bias of
Figure~\ref{fig:curves}, and the defense is the same both times:
decide the evaluation before the thing being evaluated exists. This
chapter is that decision. It collects every falsification promise
made in Chapters~\ref{sec:modelfree} through~\ref{sec:product} into
numbered tests with pass rules, prices the expensive ones, and ties
each build tier's completion to named entries rather than to the
feeling of being done.

\section{Three layers}

The suite runs at three costs, and a test's layer is part of its
identity.

Layer one is cheap and broad: surrogate and tabular benchmarks
(YAHPO-Gym, HPOBench, and the RL-specific ARLBench), where a full
tuner comparison costs CPU-minutes because the training runs were
paid for once by others \citep{pfisterer2022yahpo,
eggensperger2021hpobench, becktepe2024arlbench}, plus recorded-curve
replays in the style of Syne Tune's simulator
(Chapter~\ref{sec:libraries}), where a scheduler's decisions are
re-executed against stored learning curves in seconds. A scheduler
bug found in a CPU-minute replay is a bug not found in a GPU-week
study; layer one exists to make that sentence routine. The
multi-fidelity-with-priors suite from the PriorBand line covers the
prior-handling code specifically \citep{mallik2023priorband}.

Layer two is realistic and ours: one pinned task per workload
family. A PPO task on a Brax environment; a flow-plus-HMC sampler
task scored by ESS per gradient under divergence and
split-$\widehat{R}$ vetoes, with a reference posterior computed once
at great length; a Hamiltonian network task scored on prediction
error and energy drift. Layer two is where verdicts about \emph{our}
workloads are settled, at GPU prices, under the statistical
discipline of Section~\ref{sec:val-protocol}.

Layer three is standing: diagnostics that run inside every real
study the product ever executes. The rank correlation between rung
scores and final scores is logged always
(Chapter~\ref{sec:multifidelity}); the seed protocol's presence or
explicitly recorded absence is part of every study row
(Chapter~\ref{sec:contracts}). Layer three is how the suite keeps
learning after release, and it costs nothing because the store
records it anyway.

\section{The register of tests}

Table~\ref{tab:suite} is the register. Each entry names the claim
under test, the protocol in one line, the rule that passes or vetoes,
the chapter whose promise it discharges, and the tier gate it
belongs to. The subsections after it expand the entries whose
protocols need more than a line.

\begin{table}[tp]
\centering
\footnotesize
\renewcommand{\arraystretch}{0.95}
\caption{The validation register. Layer 1 runs on tabular benchmarks
and replays (CPU); layer 2 on the pinned tasks (GPU); layer 3 inside
every production study. ``Gate'' names the tier that cannot be
declared done until the entry is green.}
\label{tab:suite}
\begin{tabularx}{\textwidth}{@{}l>{\raggedright\arraybackslash}X>{\raggedright\arraybackslash}Xll@{}}
\toprule
ID & Claim under test & Pass rule or veto & Evid. & Gate \\
\midrule
V01 & each wrapped searcher/scheduler matches its reference
implementation on tabular benchmarks & final and anytime scores
within tolerance of reference & \ref{sec:libraries} & T0 \\
V02 & scheduler logic is correct under asynchrony & replayed
decisions on recorded curves match specification; stragglers and
drops simulated & \ref{sec:multifidelity} & T0 \\
V03 & the suite itself catches wiring bugs & seeded defects
(off-by-one rung, dropped seed, swapped objective sign) are caught
by V01--V02 & \ref{sec:product} & T0 \\
V04 & every shipped searcher earns its complexity & beats the
log+Sobol floor at matched budget on pinned tasks, each searcher
tested at the gate of the tier that ships it: the tier-0 GP with
qLogEI and the wrapped TPE at the tier-0 gate, the tier-1 and later
searchers at theirs & \ref{sec:modelfree}, \ref{sec:bayesian} &
T0/T1 \\
V05 & the transform matters as claimed & GP with and without
log-warping on a pinned task; warped must not lose & \ref{sec:modelfree} & T0 \\
V06 & early stopping is nearly free & searcher+ASHA matches
searcher-at-full-fidelity final quality within noise at a fraction
of compute; rung correlations logged & \ref{sec:multifidelity} & T1 \\
V07 & the small-budget RL verdict holds here & DEHB vs BO+ASHA at
matched 16-run budget on the pinned RL task & \ref{sec:multifidelity},
\ref{sec:population} & T3 \\
V08 & the population line does not lose in its home regime & population line vs
DEHB and BO+ASHA on both sides of the eight-worker line at matched
compute; 3 reps/arm/side (declared exception,
Sec.~\ref{sec:val-v08}); consistent home loss demotes the flagship,
disagreement reports unresolved & \ref{sec:population} & T2 \\
V09 & the premium MO searcher earns its tier & qLogNEHVI vs
random-weight Chebyshev and wrapped NSGA-II on the Hamiltonian task;
hypervolume with uncertainty; demote to option on a tie & \ref{sec:moo} & T1 \\
V10 & cheap MO fidelity signals may cull & short-horizon drift
passes the rank-correlation pilot before MO-ASHA may kill on it & \ref{sec:moo} & T1 \\
V11 & priors are safe to default-on & no-prior vs folklore-prior vs
wrong-prior runs; gains under good, recovery within noise under
wrong, else demote to opt-in & \ref{sec:transfer} & T1 \\
V12 & warm starts pay & measured trial savings on a new study of a
familiar family once the store holds enough studies & \ref{sec:transfer} & T1$^{+}$ \\
V13 & sampler correctness cannot be bought by speed & veto-failing
configurations never promoted; survivors ranked by ESS/gradient;
pilot agreement with the reference posterior & \ref{sec:workloads} & T1 \\
V14 & tier 0 is composition, not construction & the day-one walk of
Chapter~\ref{sec:product} reproduced end to end within the tier-0
budget & \ref{sec:product} & T0 \\
V15 & the curve-model bet is live before we pay for it & ifBO
matches or beats ASHA on the curve-informative public suites before
any integration work & \ref{sec:multifidelity} & T3 \\
\bottomrule
\end{tabularx}
\end{table}

\section{The statistical protocol}
\label{sec:val-protocol}

Every layer-two comparison follows one protocol, copied from the
discipline of \citet{eimer2023hyperparameters} and from this
document's own warnings. Fixed tuning seeds and disjoint test seeds.
Budgets matched on measured cost, with the tuner's own overhead
included in the bill. Results reported as performance-over-budget
curves with uncertainty bands, never as single numbers. And a
repetition count that is derived rather than declared: a pilot
campaign estimates the between-run variance, and the sample size
follows from that variance and the entry's decision margin at 80
percent power. The ten repetitions this document has quoted since
Chapter~\ref{sec:workloads} are a planning heuristic for the
calendar and the budget, not a substitute for that calculation.
Where the pilot shows low variance against a wide margin, fewer
repetitions suffice and the protocol records why; where it shows
high variance against a narrow one, the budget rises or the claim
softens, and the entry says which.

Because the budget is finite, every entry also carries a maximum
sample size, and reaching it without crossing the decision threshold
is a third outcome alongside pass and fail: \emph{inconclusive at
budgeted precision}. That is a real verdict with a real consequence,
which is that the capability does not become a default on evidence
that never arrived. Pilot variance estimates are recorded before the
full campaign begins and are not revised during it, since a variance
estimate updated mid-campaign is a decision threshold updated
mid-campaign wearing a disguise. V08 remains the one entry whose arms
are too expensive for any of this at full strength; it runs three
repetitions per arm per side under the restricted question,
asymmetric demotion rule, and release semantics of
Section~\ref{sec:val-v08}, and any future exception is written into
this section before its campaign runs, never after its results are
seen.

Verdicts then come in three strengths, and the suite's reports use
the distinction explicitly. Hard vetoes (a diverged run promoted, a
veto-failing configuration surviving, a protected seed leaking into
tuning) fail a test regardless of scores. Descriptive differences
(one curve above another within overlapping bands) are recorded and
never rank anything. A ranking is claimed only when the repetitions
support it statistically; when a result shows one method ahead
within noise, the honest summary is that they tied, and a tie
against a simpler method is a loss for the complex one wherever the
register says so (V04, V09). The suite exists to catch regressions
and to rank where gaps are real, not to manufacture leaderboard
drama.

\subsection{Register shorthand resolved}

The table above uses unquantified phrases to fit on one page; this
section pins each one to a concrete rule before any entry runs.

\textbf{``Within tolerance'' (V01):} Anytime and final scores differ
by no more than $0.5\%$ of the reference method's final value on each
benchmark, and rank orderings match exactly. The tolerance is loose
by design—V01 gates implementation correctness, not optimizer
superiority.

\textbf{``Within noise'' (V06, V11):} The 95\% bootstrap confidence
interval of the difference in means lies entirely inside the
preregistered non-inferiority margin $[-\delta, +\delta]$: for V06,
$\delta$ is 5 percent of the full-fidelity baseline's final quality;
for V11, 10 percent of the no-prior baseline's. Bootstrap by
resampling tuning runs with replacement, 10,000 iterations,
empirical percentiles. The distinction from an interval that merely
overlaps zero is the whole point. Overlapping zero says a difference
was not detected, which a small enough sample says about everything;
an interval contained in the margin says any difference that
survives is too small to care about. An interval that overlaps zero
but spills past $\pm\delta$ is recorded as inconclusive, not as a
tie, and a tie is claimed only in the contained case. Both margins
are set here, before the campaigns, as the smallest difference that
would change what the product recommends.

\textbf{``A fraction of compute'' (V06):} Early stopping reaches the
same final quality, within the non-inferiority margin above, at no
more than one-third the total compute of full-fidelity search.
Compute is the accelerator-seconds and wall-clock the executor
actually recorded, which puts startup, training, evaluation,
checkpoint save and load, queue wait, and the optimizer's own
overhead on the bill; the savings claim is about resources consumed,
not about an idealization of them. Full-run equivalents survive as a
normalized secondary measure for readability, computed by dividing
total cost by the mean cost of one full-fidelity trial in the
baseline arm, and they are never the number the gate reads. The
threshold applies to accelerator-seconds. Where wall-clock diverges
materially from it through scheduling or parallelism, both are
reported and the verdict follows the accelerator-seconds, since that
is the resource the budget buys. V06 passes only if ASHA delivers the
factor-of-three saving on measured cost.

\textbf{``Measured trial savings'' (V12):} The warm-started study
reaches the same final objective value in at least 20\% fewer trials
than a cold start on the same task. Both arms use the same searcher;
only the initialization differs. Below 20\%, warm starting is not
worth defaulting on.

\textbf{``Hypervolume with uncertainty'' (V09):} Report hypervolume
over budget with 95\% confidence bands from bootstrap resampling of
tuning runs. Rank methods by their final hypervolume; demote the
premium method (qLogNEHVI) to an option if the interval overlaps the
simpler method's (Chebyshev) or if it loses outright.

\textbf{``Enough studies'' (V12):} The run store must contain at
least ten completed studies from the same model family (same
architecture class, same task domain) before V12 runs. Below that,
the warm-start query has no data to learn from and the test is
meaningless.

All numerical gates (0.5\%, one-third, 20\%, 0.6 for V10, 1.01 for
V13) are judgments recorded before the data arrives, not fit to
results. The V10 and V13 protocols later in this chapter set their
thresholds explicitly; the phrases above inherit the same discipline.

\subsection{Sampling levels and paired designs}

A tuning comparison has more than one source of randomness, and
conflating them is the fastest way to a confident wrong answer. Five
levels are in play: the optimizer seed (which tuning run), the
workload training seed (which initialization or environment
instance), the evaluation seed (which test episodes), the task
(which benchmark), and the configuration draw (which candidates the
searcher happened to try). Treating the leaves as independent samples
when the branches are shared is pseudoreplication: it shrinks the
reported interval without adding information, and every gate
downstream inherits the false precision.

The first defense is pairing. Where a comparison admits it, the
methods share the optimizer seed set, the task and environment
instances, and any baseline configurations drawn once and evaluated
by all arms. Pairing removes the between-instance variance from the
contrast rather than averaging over it, which buys power at no extra
compute, and it applies wherever the methods being compared consume
seeds the same way. Where they do not, because a population method
spends seeds on a schedule a sequential method has no analogue for,
the entry says so and falls back to unpaired sampling with the larger
sample size the power calculation then demands.

Uncertainty is computed by hierarchical bootstrap rather than by
resampling individual trials. The resampling unit is the complete
tuning run; when runs span several tasks, tasks are resampled and
then runs within the sampled tasks, preserving the nesting. The
statistic of interest is recomputed on each of 10,000 bootstrap
samples and reported at empirical percentiles. This propagates every
level's variance without assuming a distributional form for any of
them, and it is the procedure behind every interval quoted in this
chapter.

\subsection{Development and acceptance tasks}

The trap this chapter opened with applies to the suite itself. A
pinned task that implementers run daily while debugging stops being a
held-out measurement of anything; it becomes the thing the
implementation was tuned on, and the gate it anchors measures fit to
the validation set. So each workload family carries two disjoint task
sets. Development tasks are named in week one of tier~0, recorded in
the validation manifest, and free to run as often as anyone likes.
Acceptance tasks are chosen by an engineer not implementing the
capability under test, revealed at gate time, and recorded in the
gate manifest with their data checksums.

Both numbers are reported, and the gap between them is itself a
diagnostic: a capability strong on development tasks and weak on
acceptance tasks has overfit the suite, which is a finding rather
than a mystery. Promotion decisions weigh the acceptance results,
because those are the only ones the implementation has not seen.

Task sets rotate on a version, not on a feeling. When a tier's gates
are recorded, new tasks enter development and some older development
tasks may migrate to acceptance for the next tier, with version
numbers and dates in the task registry. The registry is what makes
task-selection bias auditable after the fact, which matters
precisely because the person selecting is usually the person hoping.

\subsection{Multiple comparisons and endpoints}

Several entries ask more than one question. V04 compares searchers
across tasks and budget points; V09 crosses objectives with
comparators; V13 inspects many posterior marginals at once. Twenty
independent comparisons at the conventional threshold produce one
spurious winner on average, which is enough to promote a searcher
that does nothing.

Each entry therefore predeclares one primary endpoint, the single
comparison the tier's completion depends on, and labels the rest.
The primary endpoint is judged at $\alpha = 0.05$ uncorrected,
because there is one of it. Secondary comparisons that feed a ranking
claim are controlled family-wise at $\alpha = 0.05$ by the
Holm--Bonferroni stepdown procedure \citep{holm1979simple}, which is
uniformly more powerful than plain Bonferroni at the same guarantee
and needs no assumption about dependence among the tests.
Comparisons that rank nothing and decide nothing, the descriptive
kind this chapter has already licensed, need no correction, and
saying so is what keeps multiplicity control from becoming a tax on
honest reporting.

V04's primary endpoint, as the worked instance: whether the tier-0 GP
with qLogEI beats the log-plus-Sobol floor on the acceptance task at
matched budget. Comparisons among GP variants, or across development
tasks, are secondary and corrected. Every entry's primary endpoint is
written into its protocol before its campaign begins, for the same
reason its thresholds are.

\subsection{Immutable artifacts and preregistration}

None of the above binds unless the plan provably predates the data,
so every gate campaign emits a signed manifest before it runs. The
manifest is JSON and carries the protocol version and entry
identifier (\texttt{V04-T0-v1.2} and the like), SHA-256 checksums of
task code, data files, and environment configuration, the exact
tuning and test seed sets per method and repetition, the analysis
script with its thresholds and tests committed to version control
while the results do not yet exist, a full dependency lock and
hardware description, the commit hash of the code under test, and a
creation timestamp with a signature.

Results are written as versioned JSON with the manifest hash
embedded, under \texttt{validation/artifacts/}, a path deliberately
not ignored by \texttt{.gitignore}. Reanalysis is permitted and
expected; what it may not do is quietly replace the first analysis. A
changed threshold, a changed script, or a changed seed set after
results are visible produces a new manifest at an incremented version
with the reason written down, and both versions survive.

The point is not ceremony. Researcher degrees of freedom
\citep{simmons2011false} are invisible in a report and obvious in a
diff, and this is the cheapest available way to put them in a diff: a
reviewer can check that the analysis plan predates the data, that the
tested code matches the history, and that the tasks and seeds were
not chosen after the fact to flatter a result the team had already
grown attached to.

\section{The two expensive campaigns}
\label{sec:val-v08}

Two entries dominate the compute bill and deserve their design
stated in advance.

V07 races DEHB against BO+ASHA at a matched sixteen-full-run budget
on the pinned RL task, ten repetitions each. The cost driver is the
pinned task's full-run price, so the task is sized in week one with
ARLBench's guidance on representative subsets, which cut tuner
evaluation cost by roughly an order of magnitude in their study
\citep{becktepe2024arlbench}. As a sizing illustration, labeled an
estimate: if one full run costs two accelerator-hours on the chosen
environment, the campaign is $2 \times 16 \times 10 \times 2 = 640$
accelerator-hours, under a month on a single dedicated node and a
weekend on eight. If the chosen environment cannot fit that
envelope, the environment is wrong, not the protocol.

V08 is the flagship's trial, run late, when the population line
exists: the population method against DEHB and BO+ASHA on both sides
of the eight-worker boundary at matched total compute. It is the one
entry where the ten-repetition rule collides with the price of the
test, because a single population arm already costs eight
full-length trainings, and ten repetitions on both sides of the
boundary is a quarter-million accelerator-hour campaign we will not
buy. The exception is therefore stated as a protocol rather than
waved through. The repetition count drops to three per arm per side,
the smallest number that still yields a spread rather than a point,
and the reduction is paid for by narrowing the question: V08 does not
ask which method is best, only whether the boundary sits where
Chapter~\ref{sec:population} says it does. Three arms, two sides,
three repetitions, and the ordering of arms within a side is what the
entry reports, with the observed spread shown always.

The stopping rule is written in advance, because a cheap test with a
flexible endpoint is worse than no test. If all three repetitions on
a side agree in ordering, the side is called and no further
repetitions are bought. If they disagree, the entry reports
\emph{unresolved at this budget} and the flagship keeps whatever
status it had going in, which is a hold rather than a promotion; the
one outcome V08 may never produce is a promotion on evidence that did
not separate. The demotion direction is asymmetric and deliberately
so: a consistent loss in the flagship's home regime demotes it on
three repetitions, because the claim being falsified is a strong one
and a strong claim can be broken by a weak test, whereas confirming
the boundary would need the full protocol we declined to buy. That
asymmetry is the honest content of a three-repetition campaign, and
it is why V08 is written here as a falsifier of
Chapter~\ref{sec:population}'s threshold rather than as its
validation.

The result states and their release consequences are named too, since
a falsifier with no defined pass state cannot gate a tier. V08 opens
at \emph{hold: insufficient evidence} and can reach three places.
\emph{Demote}: all three repetitions on a side agree that the
population line loses in its home regime, in which case the flagship
becomes an experimental option and the roadmap's defaults revert to
whichever of DEHB or BO+ASHA won V07. \emph{Hold}: the repetitions
disagree, or the home-regime side is inconclusive, in which case the
entry stays at \emph{unresolved at this budget} and the flagship
keeps its tier-2 experimental status without becoming a default
recommendation. \emph{Pass}: reachable only through the full
sequential campaign this entry declined to buy, at up to ten
repetitions per arm per side with preregistered margins and decision
boundaries. The three-repetition run is a tripwire, not promotion
evidence, and no accumulation of tripwires becomes one.

Tier~2 completion therefore requires that V08 be neither demoted nor
blocked by an implementation failure; a hold is acceptable
completion. That is the honest reading of a capability whose audience
is users at eight or more dedicated workers: the population line
ships as an option with Chapter~\ref{sec:population}'s boundary
documented as guidance rather than enforced as a default, and the
campaign that would promote it waits for tier~3 or a post-release
cycle, if demand and budget ever arrive together.

\section{Two entries whose protocol needs pinning}

V10 and V13 both gate a capability on a threshold, so the threshold
has to be a number chosen before the data arrives rather than after.

V10 asks whether a short-horizon energy-drift measurement may cull a
Hamiltonian trial. The pilot runs a fixed set of one hundred
configurations drawn uniformly from the search space, each at a
single fixed seed, to full fidelity (the study's declared endpoint),
recording drift at the candidate rung (proposed as one third of the
endpoint) and at the end. The gate metric is Spearman rank
correlation between the two drift vectors (one value per
configuration, not per run, since single-seed eliminates within-config
variance). Compute the bootstrap 95\% confidence interval by
resampling configurations with replacement 10,000 times and taking
empirical percentiles of the resampled correlations. MO-ASHA may kill
on that signal only if the interval's lower bound clears $0.6$. Below
it, the scheduler still runs but the signal is advisory, promotion
decisions falling back to the objective that does correlate. The
number 0.6 is a judgment, not a result from the literature, and it is
recorded here so that a later argument about it is an argument about
a written-down choice. The hundred-trial campaign costs about 133
full runs on the Hamiltonian task, budgeted under layer-two
validation.

V13 asks whether the sampler study's correctness vetoes hold. Three
components, each with its own rule. The vetoes themselves are
absolute: any configuration with a divergent transition or
split-$\widehat{R}$ above $1.01$ is never promoted, and a single
promotion of a veto-failing configuration fails the entry outright.
The sampling protocol that feeds the vetoes: four independent chains
per configuration, each from a dispersed random initialization, 2,000
warmup draws discarded plus 2,000 retained draws per chain, for 8,000
posterior samples total per configuration. Split-$\widehat{R}$ is
computed by splitting each chain in half (eight half-chains) and
comparing within-half-chain variance to between-half-chain variance
per the rank-normalized folded formula of \citet{vehtari2021rhat}; the
veto fires if any marginal exceeds $1.01$ or if effective sample size
falls below 400. Divergent transitions are reported by the sampler
itself (a numerical-integration failure during the Hamiltonian
trajectory). A bimodal or multimodal target is detected by visual
inspection of marginal trace plots during the pilot; if found, chains
are re-initialized from dispersed modes rather than randomly, per
Bayesian workflow recommendations.

Survivors are ranked by effective sample size per gradient evaluation,
the cost measure the executor records rather than one the user
supplies. And agreement with the reference posterior is checked on the
pilot: the reference is computed once at great length using the same
four-chain protocol but with 20,000 retained draws per chain, and each
candidate configuration's marginal means and variances are compared to
the reference marginals. Agreement is defined as: the candidate's
posterior mean falls within one standard error of the reference
posterior mean, where the standard error of each mean is $\sigma /
\sqrt{\text{ESS}}$, using each sampler's own reported effective sample
size and posterior standard deviation. An entry that ranks a fast,
subtly wrong sampler above a slow, correct one is the failure this
test exists to catch, which is why speed enters only after the vetoes
and the agreement check have passed.

\section{Sabotage entries}

Two entries test the testers. V03 seeds defects into the wiring, an
off-by-one rung, a dropped seed, a flipped objective sign, and
requires the parity and replay layers to catch them; a suite that
cannot catch a planted bug is decoration, and this entry is the
reason we can trust green elsewhere. V11 seeds bad advice, the
purposefully detrimental prior in $\pi$BO's own sense, and requires
the mechanism's safety half, recovery to the prior-free baseline, to
hold on our tasks and budgets, not only in the paper
(Chapter~\ref{sec:transfer}).

\section{What the suite costs, and when it runs}

Layer one is effectively free and runs on every change: minutes of
CPU per full pass, which is why it can gate merges. Layer three is
free by construction. Layer two is a budget line: V04, V05, V06,
V09, V10, V11, and V13 together are dominated by the RL entries and
the ten-repetition rule; on the sizing illustration above, a full
tier-1 gate pass is of the order of a few hundred accelerator-hours,
scheduled as a campaign with a manifest, not run ad hoc. V07 and V08
are separate, priced campaigns. All layer-two numbers are estimates
until the week-one task sizing; the envelope, one GPU-week class of
compute per tier gate, is the planning figure the calendar of
Chapter~\ref{sec:roadmap} carries.

\section{Gates, and what failure does}

A tier is done when its register entries are green, and each entry's
failure has a named consequence rather than a shrug. Wiring failures
(V01--V03, V14) block the tier and are repaired. Evidence failures
execute the demotion already promised in the verdict's chapter:
V04 removes the losing searcher from defaults, V09 demotes the
premium searcher to an option, V11 demotes priors to opt-in, V08
demotes the flagship. The register of
\texttt{product\_register.json} records, for every verdict, which
entry anchors it; a verdict whose entry fails does not quietly
survive in the code while dying in the appendix.

\section{What would overturn this chapter}

The suite is itself under test. If V03's planted defects pass, the
parity layer is rebuilt before anything else proceeds. If the
layer-two envelope exceeds its planning figure by more than a small
factor after week-one sizing, the pinned tasks are resized and this
chapter's estimates corrected in place. And if layer three's
standing diagnostics show our real studies violating the
assumptions the pinned tasks embody, fidelities decorrelating where
the suite assumed they correlate, the pinned tasks follow the
evidence, because a suite that drifts from the workloads it guards
inherits exactly the proxy bias it was built to catch.
